\section{VI. Thảo luận và mở
rộng.}\label{vi.-thux1ea3o-luux1eadn-vuxe0-mux1edf-rux1ed9ng.}

Từ việc phân tích có thể thấy các yếu tố ảnh hưởng đến độ nhám là:
\texttt{layer\_height}, \texttt{nozzle\_temperature},
\texttt{bed\_temperature}, \texttt{material}, cuối cùng là
\texttt{print\_speed}.

\(R^2\) hiệu chỉnh = 0.8571 cho thấy mô hình ta xây dựng tương đối phù
hợp trong việc thực hiện dự báo.

Ưu điểm của phương pháp hồi quy chính là một phương pháp thống kê để
thiết lập mối quan hệ giữa một biến phụ thuộc và một nhóm tập hợp các
biến độc lập. Mô hình với một biến phụ thuộc với hai hoặc nhiều biến độc
lập được gọi là hồi quy bội (hay còn gọi là hồi quy đa biến).

Hồi quy tuyến tính bội dễ hiểu và dễ triển khai, đặc biệt khi các biến
có mối quan hệ tuyến tính và hiệu quả tính toán của phương pháp này có
thể được tính toán nhanh chóng.

Dự đoán chính xác các biến độc lập có mối quan hệ tuyến tính mạnh với
biến phụ thuộc, phương pháp hồi quy tuyến tính bội có thể cung cấp các
dự đoán chính xác, các hệ số hồi quy tuyến tính bội cung cấp thông tin
về mối quan hệ giữa từng biến độc lập và biến phụ thuộc, giúp hiểu rõ
hơn về ảnh hưởng của từng yếu tố. và tiện lợi cho kiểm định giả thuyết,
dễ dàng thực hiện các kiểm định giả thuyết về các hệ số hồi quy để kiểm
tra sự ảnh hưởng của các biến độc lập ngay cả với dữ liệu lớn, nhờ các
thuật toán tối ưu hóa hiệu quả.

Tuy nhiên mô hình cũng có một số nhược điểm như: Phương pháp này giả
định rằng mối quan hệ giữa các biến độc lập và biến phụ thuộc là tuyến
tính, điều này không phải lúc nào cũng đúng trong thực tế, hồi quy tuyến
tính bội rất nhạy cảm với các điểm dữ liệu ngoại lệ, có thể làm sai lệch
mô hình.

Khi các biến độc lập có tương quan cao với nhau, nó có thể gây ra đa
cộng tuyến, làm cho các ước lượng hệ số hồi quy không ổn định và khó
diễn giải. Phương pháp này giả định rằng các sai số có phân phối chuẩn
và có phương sai không đổi. Nếu các giả định này không được thỏa mãn,
kết quả hồi quy có thể không tin cậy.

Không phù hợp cho các mối quan hệ phi tuyến tính: Đối với các mối quan
hệ phi tuyến tính giữa các biến, hồi quy tuyến tính bội không phải là
lựa chọn tốt nhất. Các phương pháp khác như hồi quy phi tuyến hoặc mô
hình cây quyết định có thể phù hợp hơn.

Ngoài ra, đề tài ta có thể phân tích xây dựng bổ sung hai mô hình hồi
quy của hai biến \texttt{roughness} và \texttt{elongation} dựa trên cách
thức giống như xây dựng mô hình ứng với biến \texttt{tension\_strenght}.

Từ những phương pháp này mà đề tài đã đánh giá được hiệu quả những yếu
tố của biến ảnh hướng đến các sức căng bề in 3D theo như bài toán đã
được đặt ra.
\newpage